\clearpage\fancyhead[L]{\small PALLAS.}\hypertarget{native-354-lower}{}\label{native-354-lower}
\begin{center}\textbf{PALLAS.}\par\vspace{4pt}\hrule\vspace{4pt}\fontsize{10}{12}\selectfont Göttingische gelehrte Anzeigen. No. 67, pp. 657\;..\;660, 25 April 1812.\par\vspace{4pt}\hrule\vspace{4pt}\end{center}

\par\indent We have already mentioned Professor \textsc{Gauss}’s work on the perturbations of Pallas several times in these pages (see nos.~8 and 130 of last year [pp.~322 and 327 of this volume]). The results of those investigations outlined there have now been well confirmed once again. On 4 April the weather allowed a search for the planet for the first time this year. A little star of the eleventh magnitude was visible at exactly the spot indicated by the ephemeris Mr. \textsc{Nicolai} had calculated from the theory of perturbations. On 9 April it had left that spot and moved on as expected. A couple of comparisons with 93 Herculis yielded the following position, though the declination is only an estimate.

\begin{center}\fontsize{10}{12}\selectfont\setlength{\tabcolsep}{0pt}\renewcommand{\arraystretch}{1.15}\begin{tabular}{l}
April 9. $11^{\mathrm u}\;32^{\mathrm m}\;27^{\mathrm s}$ M. T.\\
Right ascension of Pallas $269^\circ\;6'\;6''7$\\
Northern declination $16\;38$
\end{tabular}\end{center}
\SourceLeaf{355}\Needspace{13\baselineskip}\par\indent This position of the planet agrees perfectly with that ephemeris. Soon afterward Professor \textsc{Gauss} heard from Dr. \textsc{Olbers}, whom he had asked to look for the planet, that he had already found it on 3 April. The two Bremen observations below likewise agree as closely as possible with Mr. \textsc{Nicolai}’s ephemeris.

\begin{center}
\fontsize{10}{12}\selectfont\setlength{\tabcolsep}{3pt}\renewcommand{\arraystretch}{1.15}\begin{tabular}{c|c|c}
Mean time & \shortstack{Right\\ascension} & \shortstack{Northern\\declination} \\
\hline
April 3 \quad $11^{\mathrm u}\;26^{\mathrm m}\;44^{\mathrm s}$ & $268^\circ\;42'\;23''1$ & $15^\circ\;25'\;4''9$ \\
\phantom{April} 4 \qquad\quad $12\;2\;33$ & $268\;47\;12.5$ & $15\;37\;41.1$ \\
\end{tabular}
\end{center}

\par\indent We hasten to make these positions known. The planet is extremely difficult to observe this year because it is so faint, and we earnestly urge all astronomers to observe it carefully, especially around the time of opposition (10 June).

\par\indent Professor \textsc{Gauss} has also recently completed the \emph{first preliminary} calculation of the \emph{general} theory of the perturbations of Pallas. The total number of periodic equations due to Jupiter’s influence that are still appreciable reaches \emph{Four hundred}. One of the very remarkable results already obtained from this work is of the greatest interest. For certain reasons we record it here in the following cipher, for which we will give the key in due course.

\begin{center}
1111000100101001
\end{center}

\par\indent We take this opportunity to add two other interesting astronomical communications, from a letter dated 26 March from Professor \textsc{Bessel} in Königsberg to Professor \textsc{Gauss}.

\par\indent As is well known, \textsc{Herschel} had already put forward the view that stars appearing to us as double stars do not line up with us by chance. They really are close to each other, and each pair forms a system of its own. Although there are still substantial objections to \emph{the reasons} that led this famous observer to draw that conclusion, other observations speak very strongly in favor of the view itself. On grounds of probability theory, it is extremely unlikely that \emph{chance} should have formed even a single pair of double stars among the stars up to a certain order, let alone as many as actually occur in the sky. This view has now received very convincing confirmation from a very interesting observation by Professor \textsc{Bessel}. His important study of \textsc{Bradley}’s observations, on which he has been working for several years, has enabled him to identify several fixed stars with very large proper motions. Most remarkable among these is the double star 61 Cygni, whose proper motion is the largest known. Both stars share this motion. It is directed toward Hercules, a direction that all large proper motions also follow very closely. For 1755, Mr. \textsc{Bessel} finds the following from \textsc{Bradley}’s observations.

\begin{center}
\fontsize{10}{12}\selectfont\setlength{\tabcolsep}{3pt}\renewcommand{\arraystretch}{1.15}\begin{tabular}{c|c|c}
 & \shortstack{Right\\ascension} & \shortstack{Northern\\declination} \\
\hline
61 Cygni & $313^\circ\;59'\;11''3$ & $37^\circ\;33'\;29''7$ \\
Sequens & $313\;59\;25.7$ & $47\;33\;45.7$ \\
\end{tabular}
\end{center}

\par\indent \textsc{Piazzi}’s data could not be used because the year of observation is unknown. But two observations in the \emph{Histoire céleste} give the annual motion of the primary star

\begin{center}\fontsize{10}{12}\selectfont\setlength{\tabcolsep}{0pt}\renewcommand{\arraystretch}{1.15}\begin{tabular}{l}
$+5''250$ in right ascension;\\$+2''321$ in declination.
\end{tabular}\end{center}
\SourceLeaf{356}\noindent If the annual parallax of this double star could be determined by observation, once the motion of the two stars around their common center of mass became known, it would be possible to deduce the ratio of the sum of their masses to the mass of our Sun. This shows how much this remarkable star deserves astronomers’ closest attention.

\Needspace{25\baselineskip}\par\indent A second communication from Professor \textsc{Bessel} concerns last year’s great comet. As early as eight months ago (see last year’s Gelehrte Anzeigen, p.~1293 [p.~329 of this volume]), we noted that there was hope of finding this comet again for a third time in Aquarius in April and the following months. There is still no reason to give up that hope. Provided that physical causes do not alter its brightness, the comet will still have one fifth as much light as it had at the beginning of January, when it could still be made out with the naked eye. The following short ephemeris, calculated by Professor \textsc{Bessel}, will be of real help in the search.

\begin{center}
\fontsize{10}{12}\selectfont\setlength{\tabcolsep}{3pt}\renewcommand{\arraystretch}{1.15}\begin{tabular}{c|c|c|c}
\shortstack{Mean Paris time\\$1812\;16^{\mathrm u}$} & \shortstack{Right\\asc.} & \shortstack{Southern\\decl.} & Brightness \\
\hline
April 12 & $335^\circ\;44'$ & $6^\circ\;52'$ & 0.0068 \\
16 & $336\;14$ & $7\;6$ & 0.0067 \\
20 & $336\;42$ & $7\;21$ & 0.0065 \\
24 & $337\;6$ & $7\;37$ & 0.0064 \\
28 & $337\;28$ & $7\;55$ & 0.0063 \\
May\ \ 2 & $337\;46$ & $8\;14$ & 0.0062 \\
6 & $338\;1$ & $8\;35$ & 0.0061 \\
10 & $338\;12$ & $8\;57$ & 0.0061 \\
14 & $338\;20$ & $9\;21\tfrac{1}{2}$ & 0.0060 \\
18 & $338\;23$ & $9\;48$ & 0.0060 \\
22 & $338\;23$ & $10\;16$ & 0.0059 \\
26 & $338\;18$ & $10\;46$ & 0.0059 \\
30 & $338\;10$ & $11\;19$ & 0.0058 \\
\end{tabular}
\end{center}
